% SPDX-License-Identifier: CC-BY-NC-ND-4.0
% Copyright 2026 Ingolf Lohmann.
\documentclass[11pt,a4paper]{article}
\usepackage[T1]{fontenc}
\usepackage[utf8]{inputenc}
\usepackage[ngerman]{babel}
\usepackage{lmodern}
\usepackage{microtype}
\usepackage{amsmath,amssymb,mathtools}
\usepackage{booktabs,array,longtable}
\usepackage{geometry}
\usepackage{xcolor}
\usepackage{hyperref}
\usepackage{enumitem}
\geometry{margin=27mm}
\hypersetup{colorlinks=true,linkcolor=blue!45!black,citecolor=blue!45!black,urlcolor=blue!45!black,pdftitle={Das Wirkungsquadrat der Planck-Skala},pdfauthor={Ingolf Lohmann}}
\setlist{nosep,leftmargin=*}
\newtheorem{definition}{Definition}
\newtheorem{satz}{Satz}
\newtheorem{korollar}{Korollar}
\newtheorem{bemerkung}{Bemerkung}
\newcommand{\lP}{\ell_{\mathrm P}}
\newcommand{\tP}{t_{\mathrm P}}
\newcommand{\mP}{m_{\mathrm P}}
\newcommand{\pP}{p_{\mathrm P}}
\newcommand{\EP}{E_{\mathrm P}}
\newcommand{\hbarred}{\hbar}

\title{\textbf{Das Wirkungsquadrat der Planck-Skala}\\[0.5em]
\large Eine relationale Normalform von Raum, Zeit, Impuls und Energie}
\author{Ingolf Lohmann\\Independent Researcher; QIK-VRT}
\date{25. Juli 2026}

\begin{document}
\maketitle

\begin{abstract}
Die Planck-Länge, Planck-Zeit, Planck-Masse, Planck-Impuls und
Planck-Energie werden gewöhnlich als Liste natürlicher Einheiten präsentiert.
Der vorliegende Beitrag ordnet dieselben Standardgrößen in eine geschlossene
relationale Normalform ein. Für jede masselose Quantenmode gilt bei konsistenter
Phasenkonvention
\[
  \frac{\bar\lambda}{\tau}=\frac{E}{p}=c,
  \qquad
  \bar\lambda p=\tau E=\hbar,
\]
wobei $\bar\lambda=1/k$ und $\tau=1/\omega$ sind. Die Gravitation wählt aus
dieser skalenfreien Familie den positiven Planck-Punkt aus. Dort gilt exakt
\[
  \frac{\lP}{\tP}=\frac{\EP}{\pP}=c,
  \qquad
  \lP\pP=\tP\EP=\hbar,
  \qquad
  \lP=\frac{\hbar}{\mP c}=\frac{G\mP}{c^2}.
\]
Die Gleichungen werden vollständig hergeleitet, ihre Eindeutigkeit unter den
angegebenen Voraussetzungen wird gezeigt und ihre Geltungsgrenze wird
explizit festgehalten. Eine begleitende Lean-4-Formalisierung prüft den
algebraischen Kern. Das Resultat ist eine Normalform bekannter Physik, keine
Behauptung einer bereits experimentell bestätigten diskreten Raumzeit und
keine Erklärung des Ursprungs des Universums.
\end{abstract}

\section{Erkenntnisstatus und Fragestellung}

Die Einzelbeziehungen dieses Artikels gehören zur etablierten Physik:
$\omega=ck$ für masselose Moden im Vakuum, $E=\hbar\omega$,
$p=\hbar k$, $E=pc$ sowie die aus $c$, $\hbar$ und $G$ gebildeten
Planck-Größen. Der eigenständige Beitrag liegt nicht in einer Umbenennung
dieser Gleichungen, sondern in ihrer simultanen Schließung zu einem
viergliedrigen Verhältnis- und Wirkungsgefüge, in der sauberen Trennung von
universeller masseloser Struktur und gravitationsselektiertem Planck-Punkt
sowie in einer maschinenprüfbaren Beweis- und Provenienzordnung.

Der zentrale Befund lässt sich in drei Zeilen zusammenfassen:
\begin{align}
 \frac{\lP}{\tP}&=\frac{\EP}{\pP}=c, \label{eq:ratio}\\
 \lP\pP&=\tP\EP=\hbar, \label{eq:products}\\
 \lP&=\frac{\hbar}{\mP c}=\frac{G\mP}{c^2}. \label{eq:anchor}
\end{align}
Dabei übernimmt $c$ die Rolle der Verhältnisbindung, $\hbar$ die Rolle der
Wirkungsbindung und $G$ die Rolle der absoluten Skalenselektion. Die Wörter
``Bindung'' und ``Selektion'' bezeichnen hier mathematische Funktionen der
Konstanten; sie machen aus den Konstanten nicht ohne Definition zusätzliche
physikalische Operatoren.

\section{Phasenkonvention und der Faktor $2\pi$}

Eine periodische Welle kann zyklusbezogen oder phasenbezogen beschrieben
werden. Zyklusbezogen gelten
\[
 E=hf,\qquad p=\frac{h}{\lambda},\qquad c=f\lambda.
\]
Phasenbezogen setzt man
\[
 \omega=2\pi f,\qquad k=\frac{2\pi}{\lambda},\qquad h=2\pi\hbar
\]
und erhält
\[
 E=\hbar\omega,\qquad p=\hbar k,\qquad \omega=ck.
\]
Für die relationale Normalform sind die reduzierten Skalen
\begin{equation}
 \bar\lambda:=\frac{1}{k}=\frac{\lambda}{2\pi},
 \qquad
 \tau:=\frac{1}{\omega}=\frac{T}{2\pi}
 \label{eq:reduced}
\end{equation}
maßgeblich. Eine Vermischung von $h$ mit $\omega$ oder von $\hbar$ mit der
vollen Wellenlänge erzeugt einen bloßen Konventionsfehler. Die nachfolgenden
Beweise verwenden durchgehend $\hbar$, $k$, $\omega$, $\bar\lambda$ und
$\tau$.

\section{Das universelle Wirkungsquadrat}

\begin{satz}[Universelles Wirkungsquadrat]
Sei eine masselose Mode im Vakuum durch $k>0$ und $\omega=ck$ mit $c>0$
gegeben. Es gelte $p=\hbar k$, $E=\hbar\omega$ und $\hbar>0$. Mit
$\bar\lambda=1/k$ und $\tau=1/\omega$ gilt
\begin{equation}
 \frac{\bar\lambda}{\tau}=\frac{E}{p}=c,
 \qquad
 \bar\lambda p=\tau E=\hbar.
 \label{eq:universal}
\end{equation}
\end{satz}

\begin{proof}
Direktes Einsetzen ergibt
\[
 \bar\lambda p=\frac{1}{k}\hbar k=\hbar,
 \qquad
 \tau E=\frac{1}{\omega}\hbar\omega=\hbar.
\]
Ferner folgt aus $\omega=ck$
\[
 \frac{\bar\lambda}{\tau}
 =\frac{1/k}{1/\omega}=\frac{\omega}{k}=c
\]
und ebenso
\[
 \frac{E}{p}=\frac{\hbar\omega}{\hbar k}=\frac{\omega}{k}=c.
\]
\end{proof}

Die universelle Lösungsmenge besitzt einen Skalenfreiheitsgrad. Wählt man
$\bar\lambda>0$, dann folgen eindeutig
\begin{equation}
 \tau=\frac{\bar\lambda}{c},\qquad
 p=\frac{\hbar}{\bar\lambda},\qquad
 E=\frac{\hbar c}{\bar\lambda}.
 \label{eq:family}
\end{equation}
Ohne eine weitere dimensionsbehaftete Konstante wird kein einzelner Punkt der
Familie ausgezeichnet. $c$ und $\hbar$ bestimmen die Relationen, aber noch
nicht die absolute Länge.

\section{Definition der Planck-Größen}

Für $c>0$, $\hbar>0$ und $G>0$ definieren wir
\begin{align}
 \lP&:=\sqrt{\frac{\hbar G}{c^3}},
 &\tP&:=\sqrt{\frac{\hbar G}{c^5}}=\frac{\lP}{c},\\
 \mP&:=\sqrt{\frac{\hbar c}{G}},
 &\pP&:=\mP c=\sqrt{\frac{\hbar c^3}{G}},\\
 \EP&:=\mP c^2=\pP c=\sqrt{\frac{\hbar c^5}{G}}.
 \label{eq:planckdefs}
\end{align}
Alle Wurzeln bezeichnen die eindeutig positive Quadratwurzel.

\begin{satz}[Planck-Schließung]
Die Größen aus \eqref{eq:planckdefs} erfüllen exakt
\eqref{eq:ratio} und \eqref{eq:products}.
\end{satz}

\begin{proof}
Für das räumlich-zeitliche Verhältnis gilt
\[
 \frac{\lP}{\tP}
 =\frac{\sqrt{\hbar G/c^3}}{\sqrt{\hbar G/c^5}}
 =\sqrt{c^2}=c.
\]
Für Energie und Impuls folgt unmittelbar
\[
 \frac{\EP}{\pP}=\frac{\mP c^2}{\mP c}=c.
\]
Die beiden Wirkungsprodukte lauten
\[
 \lP\pP
 =\sqrt{\frac{\hbar G}{c^3}}
  \sqrt{\frac{\hbar c^3}{G}}
 =\hbar
\]
und
\[
 \tP\EP
 =\sqrt{\frac{\hbar G}{c^5}}
  \sqrt{\frac{\hbar c^5}{G}}
 =\hbar.
\]
Positivität entfernt die sonst mögliche Vorzeichenmehrdeutigkeit.
\end{proof}

\section{Gravitationsselektion}

Die reduzierte Compton-Länge einer Masse $m>0$ ist
\[
 \bar\lambda_C(m)=\frac{\hbar}{mc}.
\]
Als gravitative Längenskala verwenden wir ausdrücklich
\[
 r_g(m)=\frac{Gm}{c^2}.
\]
Diese Konvention ist die Hälfte des Schwarzschild-Radius. Die Gleichsetzung
$\bar\lambda_C(m)=r_g(m)$ ergibt
\[
 \frac{\hbar}{mc}=\frac{Gm}{c^2}
 \quad\Longleftrightarrow\quad
 m^2=\frac{\hbar c}{G}.
\]
Wegen $m>0$ folgt eindeutig $m=\mP$ und damit
\[
 \bar\lambda_C(\mP)=r_g(\mP)=\lP.
\]
Dies beweist \eqref{eq:anchor}. Verwendete man stattdessen den
Schwarzschild-Radius $2Gm/c^2$, entstünde der konventionsbedingte Faktor
$\sqrt{2}$. Daher ist die verwendete gravitative Längendefinition Bestandteil
des Satzes und darf nicht stillschweigend gewechselt werden.

\begin{satz}[Eindeutige positive Gravitationsselektion]
Sei $\ell>0$ ein Punkt der universellen Familie \eqref{eq:family}. Definiere
$m:=p/c=\hbar/(\ell c)$. Erfüllt $\ell=Gm/c^2$, dann ist
$\ell=\lP$, $m=\mP$, $p=\pP$, $E=\EP$ und $\tau=\tP$.
\end{satz}

\begin{proof}
Einsetzen von $m=\hbar/(\ell c)$ in die Gravitationsbedingung liefert
\[
 \ell=\frac{G\hbar}{\ell c^3}
 \quad\Longleftrightarrow\quad
 \ell^2=\frac{\hbar G}{c^3}.
\]
Die positive Lösung ist eindeutig $\ell=\lP$. Alle übrigen Größen folgen
eindeutig aus \eqref{eq:family}.
\end{proof}

\section{Dimensionsanalytische Eindeutigkeit}

Gesucht sei eine monomiale Länge
\[
 L=c^a\hbar^bG^d.
\]
Mit
\[
 [c]=LT^{-1},\qquad [\hbar]=ML^2T^{-1},\qquad
 [G]=M^{-1}L^3T^{-2}
\]
ergeben sich die Exponentengleichungen
\[
 b-d=0,\qquad -a-b-2d=0,\qquad a+2b+3d=1.
\]
Das lineare System besitzt die eindeutige Lösung
\[
 a=-\frac32,\qquad b=\frac12,\qquad d=\frac12,
\]
also
\[
 L=\sqrt{\frac{\hbar G}{c^3}}=\lP.
\]
Analog ergeben sich Planck-Zeit und Planck-Masse. Dimensionsanalyse beweist
hier die Eindeutigkeit der monomialen natürlichen Skala. Sie beweist für sich
allein weder eine Dynamik noch die ontische Diskretheit der Raumzeit.

\section{Kovariante Einordnung}

Die ebene Quantenphase kann geschrieben werden als
\[
 \Phi=k_\mu x^\mu=\frac{p_\mu x^\mu}{\hbar},
 \qquad p^\mu=\hbar k^\mu.
\]
In einer üblichen Signatur erscheint im Exponenten die Wirkungskombination
$\mathbf p\!\cdot\!\mathbf x-Et$. Die räumliche Paarung $x p$ und die
zeitliche Paarung $tE$ sind daher keine nachträgliche Analogie. Sie sind die
Komponenten derselben relativistischen Phasenstruktur. Das Wirkungsquadrat ist
eine skalare, didaktisch transparente Projektion dieser kovarianten Beziehung
auf charakteristische positive Skalen.

\section{Maschinenprüfbarer Beweisstatus}

Die begleitende Lean-4-Datei formalisiert die Planck-Länge als positive Wurzel
von $\hbar G/c^3$ und definiert Zeit, Masse, Impuls und Energie durch die oben
angegebenen Relationen. Der Kernel prüft insbesondere:
\begin{itemize}
 \item Positivität und Quadratgleichung der Planck-Länge,
 \item beide Verhältnisidentitäten,
 \item beide Wirkungsproduktidentitäten,
 \item quantenmechanischen und gravitativen Längenanker,
 \item Eindeutigkeit des positiven gravitationsselektierten Punktes,
 \item die zusammengefasste Drei-Zeilen-Normalform.
\end{itemize}
Das Repository-Gate untersagt \texttt{sorry}, \texttt{admit}, projektinterne
\texttt{axiom}-Deklarationen und ungebundene Beweisersatzkonstanten. Ein
Axiom-Audit dokumentiert die vom Lean-/Mathlib-Unterbau verwendeten logischen
Grundlagen. Die formale Aussage bleibt ausdrücklich an ihre Definitionen und
Positivitätsprämissen gebunden.

Ein maschineller Beweis kann zeigen, dass die Gleichungen aus den Definitionen
folgen. Er misst nicht selbst $c$, $\hbar$ oder $G$ und verwandelt keine
Korrespondenzhypothese automatisch in einen Messbefund. Diese Grenze ist Teil
der Beweisarchitektur, nicht ihre Schwäche.

\section{Was an der Darstellung neu ist}

Historische Priorität für jede Einzelgleichung wird nicht beansprucht. Der
prüfbare Beitrag besteht in der folgenden Kombination:
\begin{enumerate}
 \item der expliziten Normalform mit zwei Verhältnis- und zwei
       Wirkungsinvarianten;
 \item der strikten Trennung von universeller masseloser Lösungsfamilie und
       gravitationsselektiertem Planck-Punkt;
 \item der konsistenten Phasennormierung durch $\hbar$, $1/k$ und $1/\omega$;
 \item dem Eindeutigkeitsbeweis der positiven Selektion unter angegebener
       Gravitationskonvention;
 \item der quellen-, hash-, kernel- und veröffentlichungsgebundenen
       Reproduzierbarkeit.
\end{enumerate}
Die kürzeste sachliche Lesart lautet:
\[
 \boxed{\;c\text{ bindet die Verhältnisse},\quad
        \hbar\text{ bindet die Wirkungsprodukte},\quad
        G\text{ wählt den Planck-Maßstab.}\;}
\]

\section{Was nicht folgt}

Aus den bewiesenen Identitäten folgt nicht, dass die Raumzeit aus kubischen
Zellen der Kantenlänge $\lP$ aufgebaut ist. Es folgt nicht, dass Zeit in
Schritten von $\tP$ abläuft. Es folgt keine bestimmte Topologie, kein
Nachbarschaftsgraph, keine Übergangsdynamik, keine neue Wechselwirkung und kein
bereits bestätigtes Modell der Quantengravitation. Eine diskrete Raumzeit kann
mit der Struktur vereinbar sein; eine kontinuierliche Beschreibung mit einer
operationalen Mindestauflösung kann ebenfalls vereinbar sein. Identität der
Skalenrelationen ist nicht Identität einer Ontologie.

Für eine neue dynamische Theorie wären mindestens erforderlich: elementare
Zustände, ein Wirkungs- oder Übergangsprinzip, eine kontrollierte
Lorentzstruktur, die Reproduktion von Allgemeiner Relativitätstheorie und
Quantenfeldtheorie im geeigneten Grenzfall sowie mindestens eine neue,
quantitative und falsifizierbare Vorhersage.

\section{Die weiterhin offene Ursprungsfrage}

Die Normalform beantwortet, wie fundamentale Größen innerhalb der bekannten
physikalischen Ordnung zusammenhängen. Sie beantwortet nicht, warum überhaupt
ein Universum existiert, ob ein physikalisch sinnvoller Begriff eines
``Davor'' besteht, warum Naturgesetze mathematisch formulierbar sind oder warum
$c$, $\hbar$ und $G$ gerade ihre beobachteten Werte besitzen. Die Gleichungen
sind eine Landkarte der inneren Relationen, keine Geburtsurkunde des
Universums.

\section{Schluss}

Die Planck-Größen bilden keine lose Liste extremer Zahlen. Unter den
Standarddefinitionen bilden sie ein geschlossenes, positives
Verhältnis-Wirkungs-System:
\[
 \frac{\lP}{\tP}=\frac{\EP}{\pP}=c,
 \qquad
 \lP\pP=\tP\EP=\hbar,
 \qquad
 \lP=\frac{\hbar}{\mP c}=\frac{G\mP}{c^2}.
\]
Das Resultat ist algebraisch zwingend, dimensionsanalytisch eindeutig im
angegebenen monomialen Ansatz und maschinenprüfbar formalisiert. Seine Stärke
liegt gerade in der klaren Grenze: Bewiesen ist die relationale Normalform;
offen bleiben Dynamik, empirische Quantengravitation und Kosmogonie.

\section*{Daten-, Code- und Lizenzhinweis}
Die vollständigen Quellen, Lean-Beweise, Prüfskripte, Hashbelege und
Publikationsprovenienz werden in den beiden QIK-VRT-Repositories geführt und
über eine DOI-verknüpfte Zenodo-Fassung archiviert. Artikel und nicht
ausführbare Dokumentation stehen unter CC BY-NC-ND 4.0. Prüfcode steht unter
Apache-2.0, sofern die jeweilige Datei nichts Abweichendes bestimmt.

\begin{thebibliography}{9}
\bibitem{Planck1899}
M. Planck, \emph{Über irreversible Strahlungsvorgänge}, Sitzungsberichte der
Königlich Preußischen Akademie der Wissenschaften zu Berlin, 1899.

\bibitem{CODATA2022}
E. Tiesinga et al., \emph{CODATA Recommended Values of the Fundamental
Physical Constants: 2022}, Reviews of Modern Physics / NIST reference data.

\bibitem{LohmannManuscript}
I. Lohmann, \emph{Mandelbrot, Anschlussordnung, Physik und Retrokausalität},
Zenodo, DOI: \href{https://doi.org/10.5281/zenodo.21482023}{10.5281/zenodo.21482023}, 2026.

\bibitem{LohmannFormalization}
I. Lohmann, \emph{QIK-VRT: Vollständige maschinenprüfbare Formalisierung des
formal entscheidbaren Kerns}, Zenodo, DOI:
\href{https://doi.org/10.5281/zenodo.21488116}{10.5281/zenodo.21488116}, 2026.

\bibitem{LohmannAlpha3}
I. Lohmann, \emph{QIK-VRT: Quellengebundene Lean-4-Formalisierung des
62-seitigen Manuskripts -- Alpha 3}, Zenodo, DOI:
\href{https://doi.org/10.5281/zenodo.21529081}{10.5281/zenodo.21529081}, 2026.
\end{thebibliography}

\vfill
\noindent\textit{q.e.d.}\hfill Ingolf Lohmann
\end{document}
